\Needspace{22\baselineskip}
\subsection{Divisorial products and zero volume}
\label{subsec:zero-volume-product-obstacles}

The simplest case of the fibration reduction is a product of a curve $C$ with a compact K\"ahler manifold $Z$, with the poles of $\Psi$ along $C\times D$ for a smooth divisor $D\subset Z$. If the divisor potential on $Z$ has supremum zero, the induced zero obstacle on $C$ is zero, so the envelopes of $\Psi+w\circ p$ are determined by those of $w$ on the curve. Projective products then give potentials without analytic ordinary type whose two envelopes agree, as well as distinct models with the same zero volumes. Since top-degree mass vanishes on all these models, volumes cannot distinguish them. Polar curvature on the curve can, and it gives a complete metric on the resulting class of product models.

Let $C$ be a compact connected curve with a smooth positive form $\eta$, and let $Z$ be a connected compact K\"ahler manifold of dimension $k\geq1$. Suppose that $D\subset Z$ is a nonempty smooth reduced divisor, $\kappa$ is a K\"ahler form on $Z$, and $\psi\in\PSH(Z,\kappa)$ is smooth off $D$ and satisfies
\[
 \kappa+\ddc\psi=[D],\qquad \sup_Z\psi=0.
\]
Thus, near $D$, $\psi$ is $\log|z_1|^2$ plus a smooth function, where $z_1$ is a local equation of $D$. On $M=C\times Z$, write $p\colon M\to C$ and $r\colon M\to Z$ for the projections, and put
\[
 \Theta=p^*\eta+r^*\kappa,\qquad
 \Psi=\psi\circ r,\qquad H=C\times D.
\]
For $w\in\PSH(C,\eta)$ set $u=\Psi+w\circ p$.

On the curve $(C,\eta)$ we use the models of \cref{sec:curve-zero-volume-models}. The polar curvature of $q\in\PSH(C,\eta)$ is the restriction of $\eta+\ddc q$ to the pole set $\{q=-\infty\}$ (\cref{sec:local-polar-parts}). By \cref{lem:curve-polar-model}, every positive measure of mass at most $\int_C\eta$ carried by a polar set is the polar curvature of exactly one ordinary model, and $P_\eta[q]$ is the model with the polar curvature of $q$. Ideal projection takes the atomic part of the polar curvature (\cref{lem:curve-atomic-ideal-model,prop:curve-analytic-distance-defect}), and the total variation of the difference of polar curvatures is a complete geodesic metric on ordinary models, equal to the volume comparison $d_\eta$ of \cref{def:model-volume-comparison} (\cref{thm:curve-model-volume-distance,cor:curve-model-completeness}).

\Needspace{14\baselineskip}
\begin{proposition}\label{prop:divisorial-product-obstacles}\index{Envelope!divisorial product}\index{Singularity volume!zero-volume product}
For $F\in C^\infty(M,\mathbb R)$, the function
\[
 \widehat F(y)=\min_{z\in Z\setminus D}\{F(y,z)-\psi(z)\}
 \qquad(y\in C)
\]
is finite and continuous. The two envelopes satisfy
\begin{equation}\label{eq:divisorial-product-obstacles}
 \begin{split}
 P_{\Theta,F}[u]&=\Psi+P_{\eta,\widehat F}[w]\circ p,\\
 P_{\Theta,F}[u]_{\mathcal I}
     &=\Psi+P_{\eta,\widehat F}[w]_{\mathcal I}\circ p.
 \end{split}
\end{equation}
\end{proposition}
\begin{proof}
Apply \cref{prop:fibration-obstacle-reduction} to $p\colon M\to C$, with $H=C\times D$, $\beta=\eta$ and $\Psi=\psi\circ r$. The local logarithmic form of $\psi$ defines a smooth metric on $\mathcal O_M(H)$ with curvature $r^*\kappa$. Thus $\Theta+\ddc\Psi=p^*\eta+[H]$, and $H$ is smooth and transverse to the fibers. The induced fiber minimum is the function $\widehat F$ above. For $F=0$, the normalization $\sup_Z\psi=0$ gives $\widehat 0=0$.
\end{proof}

\Needspace{10\baselineskip}
\begin{example}\label{ex:zero-volume-nonanalytic-product}
Take $C=\mathbb P^1$ and $Z=\mathbb P^k$, $k\geq1$, with their probability Fubini--Study forms, and let $D$ be a hyperplane. In this product setting there is a potential $u$ without an analytic ordinary type such that
\[
 P_\Theta[u]=P_\Theta[u]_{\mathcal I}=\Psi,
 \qquad u\not\sim\Psi,
\]
and the two envelopes of $u$ agree below every smooth obstacle. Their common model for the zero obstacle has zero singularity volume.
\end{example}
\begin{proof}
On the curve $C$, let $w$ be the unbounded potential constructed in the proof of \cref{prop:ordinary-topology-nonseparation}, whose two envelopes are zero, and let $\psi$ be a logarithmic potential of the hyperplane $D$ with supremum zero. For $u=\Psi+w\circ p$, \cref{prop:divisorial-product-obstacles} gives the envelope identities, \cref{cor:fibration-obstacle-equality} gives equality below every smooth obstacle, and \cref{prop:fibration-obstacle-reduction} gives zero volumes. As in the proof of \cref{ex:zero-volume-unbounded-mass-fibration}, restriction to a section outside $H$ shows $u\not\sim\Psi$, and \cref{prop:analytic-envelope-type} excludes an analytic ordinary type.
\end{proof}

Vanishing of the volumes does not force the two envelopes to agree.

\Needspace{18\baselineskip}
\begin{example}\label{ex:zero-volume-distinct-product-models}\index{Envelope!distinct zero-volume models}\index{Singularity volume!distinct zero-volume models}
In the product setting, take $C=\mathbb P^1$ and $Z=\mathbb P^k$, $k\geq1$, with probability Fubini--Study forms and a hyperplane $D$. The positive closed background $\Theta$ admits normalized ordinary models $a\leq b$ with
\[
 \begin{gathered}
 P_\Theta[a]=a,\qquad
 P_\Theta[a]_{\mathcal I}=P_\Theta[b]=P_\Theta[b]_{\mathcal I}=b,\\
 a\sim_{\mathcal I}b,\qquad a\not\sim b,\\
 \underline{\operatorname{vol}}_\Theta(a)
 =\overline{\operatorname{vol}}_\Theta(a)
 =\underline{\operatorname{vol}}_\Theta(b)
 =\overline{\operatorname{vol}}_\Theta(b)=0.
 \end{gathered}
\]
Thus equal singularity volumes of a type and of its ideal model do not imply equality of the envelopes at zero volume.
\end{example}
\begin{proof}
Use the normalized diffuse polar curve model $w$ of \cref{ex:diffuse-polar-zero-model} on $\mathbb P^1$ with its probability Fubini--Study form. It satisfies $P_\eta[w]=w$ and $P_\eta[w]_{\mathcal I}=0$, and it is unbounded. Take the normalized hyperplane potential on $\mathbb P^k$ and put
\[
 a=\Psi+w\circ p,\qquad b=\Psi.
\]
The identities for the zero obstacle in \cref{prop:divisorial-product-obstacles} give all four envelope identities. Both functions have supremum zero, and $w\leq0$ gives $a\leq b$. Ideal projection preserves the ideal type by \cref{thm:hermitian-i-projection}, so $a\sim_{\mathcal I}b$. Restriction to $\mathbb P^1\times\{z\}$, for any $z\notin D$, gives the unbounded difference $w$, so $a\not\sim b$. Finally, \cref{prop:fibration-obstacle-reduction} gives zero extended mass for every ordinary representative of either model, which proves all four volume identities.
\end{proof}

A semipositive term in the background that vanishes on the slices of a product does not change the envelopes of potentials pulled back from the other factor.

\Needspace{16\baselineskip}
\begin{lemma}\label{lem:auxiliary-slice-envelopes}\index{Envelope!auxiliary product factor}
Let $Y$ and $T$ be connected compact complex manifolds, let $\pi\colon Y\times T\to Y$ be the projection, let $\gamma$ be a smooth real $(1,1)$-form on $Y$, and let $\alpha$ be a smooth semipositive real $(1,1)$-form on $Y\times T$ whose restriction to every slice $Y\times\{t\}$ is zero. Put $\Gamma=\pi^*\gamma+\alpha$. For every $u\in\PSH(Y,\gamma)$, one has $u\circ\pi\in\PSH(Y\times T,\Gamma)$ and
\[
 P_\Gamma[u\circ\pi]=P_\gamma[u]\circ\pi.
\]
If $P_\gamma[u]_{\mathcal I}$ is ordinarily equivalent to a quasi-plurisubharmonic function $g$ with $\ddc g=S-\beta'$, where $S$ is a divisorial current on $Y$, possibly zero, as in \cref{lem:divisorial-potential-domination} and $\beta'$ is a smooth form, then also
\[
 P_\Gamma[u\circ\pi]_{\mathcal I}=P_\gamma[u]_{\mathcal I}\circ\pi.
\]
\end{lemma}
\begin{proof}
Since $\alpha\geq0$, pullback maps $\PSH(Y,\gamma)$ into $\PSH(Y\times T,\Gamma)$. It preserves the bound $v\leq0$, the ordinary order and, by \cref{lem:submersion-multiplier-pullback}, multiplier-ideal equivalence. Pullbacks of competitors are therefore competitors, and since upper regularization commutes with pullback by a projection, both envelopes upstairs are at least the pullbacks of the envelopes on $Y$.

Let $v$ be an ordinary competitor for $u\circ\pi$, and replace it by the competitor $\max(v,u\circ\pi-\sup_Yu)$. For every $t\in T$, the restriction $v(\cdot,t)$ is $\gamma$-plurisubharmonic, because $\Gamma$ restricts to $\gamma$ on $Y\times\{t\}$, and it is not identically $-\infty$, nonpositive and at least as singular as $u$. Hence $v(\cdot,t)\leq P_\gamma[u]$, which proves the ordinary identity.

Under the additional hypothesis, put $q=P_\gamma[u]_{\mathcal I}$, so that $u\sim_{\mathcal I}q\sim_{\mathcal I}g$ by \cref{thm:hermitian-i-projection}. The complex Hessian of $g\circ\pi$ is $\pi^*S-\pi^*\beta'$, and $\pi^*S$ is the sum of the prime divisors $D_j\times T$ with the coefficients of the components $D_j$ of $S$; it has finite mass for a product Hermitian metric. An ideal competitor $v$ for $u\circ\pi$ satisfies $v\sim_{\mathcal I}g\circ\pi$, so \cref{lem:divisorial-potential-domination} gives $v\preceq g\circ\pi\sim q\circ\pi$. Thus $v$ is an ordinary competitor for $q\circ\pi$, and the ordinary identity and \cref{prop:p-below-i} give $v\leq P_\gamma[q]\circ\pi=q\circ\pi$.
\end{proof}

The failure of \cref{ex:zero-volume-distinct-product-models} also occurs in a nonclosed background once a second curve factor is added.

\Needspace{16\baselineskip}
\begin{example}\label{ex:nonclosed-zero-volume-distinct-models}\index{Envelope!nonclosed zero-volume obstruction}\index{Singularity volume!nonclosed zero-volume models}
In every complex dimension $n\geq3$, there is a projective manifold with the bounded mass property that admits a smooth positive nonclosed background $\Theta$ and normalized ordinary models $a\leq b$ such that
\[
 \begin{gathered}
 P_\Theta[a]=a,\qquad
 P_\Theta[a]_{\mathcal I}=P_\Theta[b]=P_\Theta[b]_{\mathcal I}=b,\\
 a\sim_{\mathcal I}b,\qquad a\not\sim b,\\
 \underline{\operatorname{vol}}_\Theta(a)
 =\overline{\operatorname{vol}}_\Theta(a)
 =\underline{\operatorname{vol}}_\Theta(b)
 =\overline{\operatorname{vol}}_\Theta(b)=0.
 \end{gathered}
\]
In particular, equality of the upper volumes of a type and of its ideal model does not force equality of the envelopes when the ideal model has zero lower volume.
\end{example}
\begin{proof}
Let $Y=\mathbb P^1\times\mathbb P^1$, with projections $\pi_1,\pi_2$, and let $\sigma$ be the probability Fubini--Study form on $\mathbb P^1$. Choose a positive nonconstant smooth function $f$ on the first factor and put
\[
 \beta=\pi_1^*\sigma+(f\circ\pi_1)\pi_2^*\sigma.
\]
This form is positive, and $\mathrm d\beta=\pi_1^*(\mathrm df)\wedge\pi_2^*\sigma\ne0$. Let $w_0$ be the normalized diffuse polar model of \cref{ex:diffuse-polar-zero-model}, which satisfies $P_\sigma[w_0]=w_0$ and $P_\sigma[w_0]_{\mathcal I}=0$, and put $w=w_0\circ\pi_1$. \Cref{lem:auxiliary-slice-envelopes}, applied to the first factor with $T=\mathbb P^1$, $\alpha=(f\circ\pi_1)\pi_2^*\sigma$ and $g=0$, gives
\[
 P_\beta[w]=w,\qquad P_\beta[w]_{\mathcal I}=0.
\]

Set $k=n-2$ and $M=Y\times\mathbb P^k$, with projections $p,r$. On the second factor take the probability Fubini--Study form $\kappa$, a hyperplane $D$, and its normalized potential $\psi$ with $\kappa+\ddc\psi=[D]$. Put
\[
 \Theta=p^*\beta+r^*\kappa,\qquad
 \Psi=\psi\circ r,\qquad a=\Psi+w\circ p,\qquad b=\Psi.
\]
The form $\Theta$ is positive and not closed. A sufficiently large multiple of the closed positive product form dominates $\Theta$, so the bounded mass property follows from Stokes' theorem for that closed form and the metric independence in \cref{prop:bmp-metric}.

Apply \cref{prop:fibration-obstacle-reduction} to $p\colon M\to Y$ and the divisor $Y\times D$, whose induced zero obstacle is zero because $\sup\psi=0$. The two identities on the base give the four identities upstairs. The remaining assertions follow as in the proof of \cref{ex:zero-volume-distinct-product-models}; for the volumes, use $\dim Y=2<n$.
\end{proof}

We return to the product $M=C\times Z$ and its fixed divisor potential $\Psi$. All the models $\Psi+q\circ p$ have zero top-degree volumes, but the polar curvature of $q$ on $C$ still distinguishes them.

\begin{definition}\label{def:divisorial-product-model-locus}\index{Model potential!divisorial product locus}
For a normalized ordinary model $q$ in background $\eta$ on $C$, put
\[
 \Phi(q)=\Psi+q\circ p,\qquad
 \mathcal S_\Psi=\{\Phi(q):P_\eta[q]=q\}.
\]
Write $\mu_q=\mathbf1_{\{q=-\infty\}}(\eta+\ddc q)$ for the polar curvature of $q$, and define
\begin{equation}\label{eq:divisorial-product-polar-metric}
 \rho_\Psi(\Phi(q_1),\Phi(q_2))
       =\|\mu_{q_1}-\mu_{q_2}\|_{\mathrm{TV}}.
\end{equation}
\end{definition}

\Needspace{15\baselineskip}
\begin{proposition}\label{prop:divisorial-product-polar-metric}\index{Metric completeness!zero-volume product}\index{Model potential!polar product metric}
The set $\mathcal S_\Psi$ consists of normalized ordinary models in background $\Theta$, and $(\mathcal S_\Psi,\rho_\Psi)$ is a complete geodesic metric space. The upper-volume comparison $d_\Theta$, given by the formula of \cref{def:model-volume-comparison} from the upper singularity volumes of $a$, $b$ and $\max(a,b)$, which are finite here, vanishes identically on this set:
\[
 d_\Theta(a,b)=0\qquad(a,b\in\mathcal S_\Psi).
\]
The ideal models in $\mathcal S_\Psi$ form a closed geodesically convex subspace, equal to the $\rho_\Psi$-closure of the models with analytic singularities. Ideal projection is a $1$-Lipschitz retraction onto this subspace and selects the unique nearest ideal model. For $a=\Phi(q)$, the distance from $a$ to this subspace is the mass of the nonatomic part of $\mu_q$. If $\rho_\Psi(a_j,a)\to0$, then $a_j\to a$ in $L^1(M)$ for every smooth positive volume form.
\end{proposition}
\begin{proof}
The identities for the zero obstacle in \cref{prop:divisorial-product-obstacles} give
\[
 P_\Theta[\Phi(q)]=\Phi(q),\qquad
 P_\Theta[\Phi(q)]_{\mathcal I}
       =\Phi(P_\eta[q]_{\mathcal I}).
\]
Both summands of $\Phi(q)$ have supremum zero on their respective factors, so $\Phi(q)$ is normalized. Restriction to any section outside $D$ shows that $\Phi$ is injective. The classification in \cref{lem:curve-polar-model} and the distance formula in \cref{thm:curve-model-volume-distance} therefore make $\Phi$ an isometry from the space of curve models onto $\mathcal S_\Psi$. Completeness and the existence of geodesics follow from \cref{cor:curve-model-completeness}.

By \cref{prop:fibration-obstacle-reduction}, every ordinary representative of $\Phi(q)$ has zero extended mass. Moreover,
\[
 \max(\Phi(q_1),\Phi(q_2))
       =\Psi+\max(q_1,q_2)\circ p,
\]
and the same proposition, applied to the curve potential $\max(q_1,q_2)$, gives zero upper volume for the maximum. Thus all three upper volumes in the definition of $d_\Theta$ vanish.

Under $\Phi$, ideal projection corresponds to taking the atomic part of $\mu_q$. The assertions on closedness, convexity and nearest points therefore follow from \cref{cor:curve-model-completeness,prop:curve-analytic-distance-defect}. To prove analytic closure upstairs, approximate an atomic measure by finite atomic measures with rational coefficients and mass less than $\int_C\eta$. By \cref{lem:curve-atomic-ideal-model}, their curve models have the corresponding logarithmic poles up to bounded remainders. Adding $\Psi$ preserves an analytic presentation: in the local product monomial, clear the denominators of the rational pole coefficients together with the unit coefficient of $D$. The resulting models in $\mathcal S_\Psi$ converge for $\rho_\Psi$. Conversely, by \cref{prop:analytic-envelope-type}, every ordinary model with analytic singularities in this locus is equal to its ideal projection, so the closure of these models lies in the closed ideal subspace.

Finally, if $\rho_\Psi(\Phi(q_j),\Phi(q))\to0$, then \cref{thm:curve-model-capacity-convergence} gives $q_j\to q$ in $L^1(C)$. Off the pole sets, which have smooth volume zero, the difference of the two product potentials is $(q_j-q)\circ p$. For a smooth positive volume form $\mathrm dV$ on $M$, integration along the compact fibers makes $p_*\mathrm dV$ a smooth positive density on $C$, and Fubini's theorem gives
\[
 \int_M|\Phi(q_j)-\Phi(q)|\,\mathrm dV
       \leq C_{\mathrm dV}\|q_j-q\|_{L^1(C,\eta)}\longrightarrow0,
\]
where $C_{\mathrm dV}$ bounds the density of $p_*\mathrm dV$ with respect to $\eta$.
\end{proof}

The locus and its metric persist when an auxiliary curve factor with a variable nonnegative coefficient is added to the background, which can make it nonclosed or nonbig; see \cref{cor:nonclosed-separated-product-metric} with $m=1$.

\Needspace{18\baselineskip}
\subsection{Several curve factors}
\label{subsec:separated-curve-product-models}

For several independent curve factors, the sum of their polar distances gives a complete metric on the induced product models. The envelope identities follow from taking a supremum over one group of factors and then restricting to its slices.

Let $C_1,\ldots,C_m$ be connected compact curves with smooth positive forms $\eta_i$, where $m\geq2$, and put $A_i=\int_{C_i}\eta_i$. Set
\[
 Y=\prod_{i=1}^m C_i,\qquad
 \beta=\sum_{i=1}^m\pi_i^*\eta_i,
\]
where $\pi_i:Y\to C_i$ are the projections. For normalized ordinary models $q_i$ on $(C_i,\eta_i)$, write
\[
 w=\sum_{i=1}^m q_i\circ\pi_i,\qquad
 q_{i,\mathcal I}=P_{\eta_i}[q_i]_{\mathcal I}.
\]

\Needspace{13\baselineskip}
\begin{lemma}\label{lem:separated-curve-model-projections}\index{Envelope!separated curve models}
The function $w$ is a normalized ordinary model on $(Y,\beta)$, and
\[
 P_\beta[w]_{\mathcal I}
       =\sum_{i=1}^m q_{i,\mathcal I}\circ\pi_i.
\]
\end{lemma}
\begin{proof}
The function $w$ is $\beta$-plurisubharmonic and has supremum zero. We prove $P_\beta[w]=w$ by induction on $m$, starting with the curve model identity. Let $v\leq0$ be an ordinary competitor, and replace it by $\max(v,w)$. This preserves $v\preceq w$ and makes the slices genuine whenever the corresponding curve potentials are finite.

Write $Y=Y'\times C_m$, and take the regularized partial supremum
\[
 a(t)=\left(\sup_{y'\in Y'}v(y',t)\right)^*.
\]
The restrictions to $C_m$ are $\eta_m$-plurisubharmonic or identically $-\infty$. They are bounded above by zero, and at least one is genuine, since $v\geq w$. The local upper-envelope theorem from \cref{chap:framework} makes $a$ $\eta_m$-plurisubharmonic. The global inequality $v\leq w+C$ gives $a\leq q_m+C$, because the other normalized potentials have supremum zero. Thus $a\leq P_{\eta_m}[q_m]=q_m$.

For $t$ where $q_m(t)$ is finite, the function $v(\cdot,t)-q_m(t)$ is a genuine potential in the product background on $Y'$, lies below zero and is more singular than $\sum_{i<m}q_i\circ\pi_i$. The induction hypothesis bounds it by that sum. Hence $v\leq w$ off the poles of $q_m$. At those poles the original bound $v\leq w+C$ gives the same inequality. Taking the envelope proves the ordinary identity.

For the ideal identity, let $E_i$ be the atomic part of the polar curvature of $q_i$, and take the Green potentials $g_i$ from \cref{lem:curve-atomic-ideal-model}. They satisfy
\[
 q_i\sim_{\mathcal I}g_i,\qquad
 q_{i,\mathcal I}\sim g_i,\qquad
 \ddc g_i=E_i-\frac{E_i(C_i)}{A_i}\eta_i.
\]
By \cref{lem:submersion-multiplier-pullback,lem:hermitian-i-calculus}, $g=\sum_i g_i\circ\pi_i$ is ideally equivalent to $w$. Its curvature is a closed countable sum of the divisors $\pi_i^{-1}(x)$ minus the smooth form $\sum_i E_i(C_i)\pi_i^*\eta_i/A_i$. Each coordinate fiber has constant mass for the product metric, so these divisor sums have finite mass. Also $g$ is $\beta$-plurisubharmonic since $E_i(C_i)\leq A_i$. Therefore \cref{cor:divisorial-envelope-equality} gives $P_\beta[g]_{\mathcal I}=P_\beta[g]$. Ordinary equivalence with $\sum_iq_{i,\mathcal I}\circ\pi_i$ and the proved ordinary identity compute this last envelope, proving the assertion.
\end{proof}

Let $Z$ be a connected compact K\"ahler manifold of dimension $k\geq1$, and let $D\subset Z$ be a nonempty smooth reduced divisor with a K\"ahler form $\kappa$ and a normalized potential $\psi$ satisfying $\kappa+\ddc\psi=[D]$. On $M=Y\times Z$, write $p,r$ for the projections, set $h_i=\pi_i\circ p$, and put
\[
 \Theta=p^*\beta+r^*\kappa,\qquad
 \Psi=\psi\circ r,\qquad
 \Phi(q_1,\ldots,q_m)=\Psi+\sum_i q_i\circ h_i.
\]
Let $\mathcal S_\Psi^{(m)}$ be the set of these potentials. For $\mu_i=\mathbf1_{\{q_i=-\infty\}}(\eta_i+\ddc q_i)$ and the analogous $\nu_i$ from models $b_i$, define
\begin{equation}\label{eq:separated-product-polar-metric}
 \rho_\Psi^{(m)}(\Phi(q_1,\ldots,q_m),\Phi(b_1,\ldots,b_m))
       =\sum_i\|\mu_i-\nu_i\|_{\mathrm{TV}}.
\end{equation}

\Needspace{15\baselineskip}
\begin{proposition}\label{prop:separated-product-polar-metric}\index{Metric completeness!several curve factors}\index{Model potential!separated product metric}
The set $\mathcal S_\Psi^{(m)}$ consists of normalized ordinary models and is a complete geodesic metric space for $\rho_\Psi^{(m)}$. Its ideal models form a closed geodesically convex subspace, equal to the metric closure of its analytic models. Ideal projection is the unique nearest-point $1$-Lipschitz retraction, and the distance to that subspace is the sum of the nonatomic polar masses. Metric convergence implies $L^1(M)$ convergence for every smooth positive volume form. The top-volume comparison $d_\Theta$ vanishes on every pair in this locus.
\end{proposition}
\begin{proof}
The preceding lemma and \cref{prop:fibration-obstacle-reduction} give
\[
 \begin{gathered}
 P_\Theta[\Phi(q_1,\ldots,q_m)]=\Phi(q_1,\ldots,q_m),\\
 P_\Theta[\Phi(q_1,\ldots,q_m)]_{\mathcal I}
       =\Phi(q_{1,\mathcal I},\ldots,q_{m,\mathcal I}).
 \end{gathered}
\]
The induced zero obstacle is zero, and independence of the factors gives normalization. Restriction to sections where the other potentials are finite shows that equality of two product potentials forces each coordinate difference to be constant. Normalization makes every such constant zero. Thus $\Phi$ identifies this locus with the finite product of the curve model spaces, equipped with the sum of their metrics.

Completeness follows coordinate by coordinate from \cref{cor:curve-model-completeness}. Interpolating each polar measure with the same parameter gives a constant-speed geodesic for the sum metric. The ideal projection takes the atomic part in each coordinate. The closedness, convexity, unique nearest-point and $1$-Lipschitz assertions follow coordinate by coordinate from \cref{cor:curve-model-completeness,prop:curve-analytic-distance-defect}; their distances add. Finite rational atomic approximations in each coordinate have local logarithmic coefficients that, together with the unit coefficient of $D$, clear denominators in a product monomial. They give analytic density in this subspace. Conversely, every analytic ordinary model is its ideal projection by \cref{prop:analytic-envelope-type}, proving that its closure is exactly the ideal subspace.

Metric convergence gives $L^1$ convergence on each curve by \cref{thm:curve-model-capacity-convergence}. The difference of the product potentials off their pole sets is the sum of the pulled-back coordinate differences. Integrating each term along the compact fibers of $h_i$ bounds its upstairs $L^1$ norm by a constant times its curve $L^1$ norm, as in \cref{prop:divisorial-product-polar-metric}. Their finite sum tends to zero.

Finally, every model in this locus has zero upper singularity volume by \cref{prop:fibration-obstacle-reduction}, since $\dim Y=m<\dim M$. A pairwise maximum is still $\Psi$ plus the pullback of a $\beta$-plurisubharmonic function on $Y$, so the same proposition gives its zero upper volume. All three terms of $d_\Theta$ therefore vanish.
\end{proof}

